\DocumentMetadata{lang=es-DO,testphase={phase-III,math}}
\documentclass[11pt]{scrartcl}
\usepackage{icv}

% -----------------------------------------------------------------------
% Migrado de legacy/Acueductos y Alcantarillados/Assignments/
% Proyección-poblacional/asignacion_proyeccion_poblacional.tex --
% asignación real completa, sin recortar. No usa \icvproblem porque no es
% una lista de preguntas puntuadas -- es un caso de estudio de varias
% partes con rúbrica analítica embebida (tabla propia, no icvrubric: el
% formato de icvrubric es para documentos de rúbrica independientes con
% los 4 niveles fijos, no para una tabla criterio/descripción/puntos
% embebida dentro de otro documento).
% -----------------------------------------------------------------------
\icvsetup{
  asignatura  = {ICV 442-T -- Acueductos y Alcantarillados},
  titulocorto = {Proyección poblacional},
  titulo      = {Asignación comprensiva de análisis y proyección poblacional},
  subtitulo   = {},
  profesor    = {Dr. Ricardo Hernández Moreira},
  periodo     = \icvciclo{2026}{3},
  version     = {v1},
}

\begin{document}
\icvmaketitle

\begin{icvnote}[Propósito]
Aplicar criterios censales, normativos y de proyección poblacional para
sustentar una población de diseño en sistemas de agua potable y
alcantarillado.
\end{icvnote}

\begin{icvnote}[Qué se espera]
Organizar datos en Excel, aplicar varios métodos, comparar resultados
con criterio ingenieril y justificar técnicamente una población de
proyecto.
\end{icvnote}

\section{Naturaleza de la actividad}
\begin{itemize}
  \item Modalidad: grupal.
  \item Producto central: portafolio técnico con soporte numérico en
    Excel y síntesis escrita breve.
  \item Referencia normativa principal: Reglamento para el Diseño de
    Obras e Instalaciones Hidro-Sanitarias del INAPA, apartado de
    estimación de población.
  \item Horizonte de proyección: cuando el enunciado pida una proyección
    futura, la población deberá estimarse para el año \textbf{2046},
    es decir, \textbf{20 años a partir de 2026}.
\end{itemize}

\section{Objetivo}
El objetivo de esta asignación es que cada grupo desarrolle criterio
para estimar y justificar la población de proyecto a partir de datos
censales, información oficial y métodos clásicos de proyección. La
actividad combina un caso con información limitada, una serie sintética
con tendencia controlada y una serie real de larga duración.

\section{Marco de referencia}
En el Reglamento del INAPA, la población futura debe estimarse a partir
de datos censales e información local o regional, y el método
seleccionado debe justificarse en función del comportamiento observado.
En consecuencia, en esta asignación interesa tanto el resultado numérico
como la calidad de la justificación.

\begin{icvwarning}[Entrega]
Todos los cálculos deben elaborarse en Excel. El informe no sustituye la
hoja de cálculo; la resume y la interpreta.
\end{icvwarning}

\section{Parte A. Caso guiado: La Siembra, Distrito Municipal Padre Las Casas, Azua}

\subsection{Contexto}
La Siembra es una comunidad con población reducida y una historia
administrativa que dificulta construir una serie larga y limpia. Se
utiliza como caso de trabajo porque obliga a verificar fuentes,
reconocer límites de los datos y justificar con prudencia la población
de diseño.

\subsection{Datos y verificación}
Cada grupo deberá identificar y citar directamente en la ONE la
población correspondiente al Censo 2010 y al Censo 2022 para La
Siembra. Si encuentra discrepancias entre tablas, publicaciones o bases
desagregadas, deberá documentarlas y explicar cuál valor adoptó.

\subsection{Actividades}
\begin{enumerate}
  \item Organice en Excel los datos censales verificados, con años,
    población y fuente.
  \item Calcule incremento absoluto, incremento porcentual y una tasa
    media representativa para el período 2010--2022.
  \item Aplique al menos dos métodos de proyección entre los
    siguientes:
    \begin{itemize}
      \item aritmético,
      \item geométrico,
      \item ajuste lineal por mínimos cuadrados,
      \item ajuste exponencial por mínimos cuadrados, si considera
        que el comportamiento observado lo justifica.
    \end{itemize}
  \item Proyecte la población al año \textbf{2046}.
  \item Analice críticamente los resultados considerando escasez de
    datos, escala poblacional, posible migración, estancamiento o
    decrecimiento y pertinencia de los métodos seleccionados.
  \item Seleccione una población de proyecto y justifíquela conforme al
    Reglamento del INAPA.
\end{enumerate}

\section{Parte B. Ejercicio con serie sintética}

\subsection{Contexto}
Esta parte utiliza una serie poblacional sintética con fines
pedagógicos. La serie fue construida para representar un patrón
compatible con crecimiento logístico, pero con variaciones pequeñas que
evitan un ajuste perfecto.

\subsection{Datos}
Para esta parte, el valor de \textbf{0 años antes del presente}
corresponde al año \textbf{2026}.

\begin{center}
\begin{tabular}{l *{5}{S[table-format=4.0]}}
  \toprule
  {Años antes del presente} & {0} & {5} & {20} & {30} & {50} \\
  \midrule
  {Población (hab)} & 7810 & 7420 & 5520 & 3460 & 1040 \\
  \bottomrule
\end{tabular}
\end{center}

\subsection{Actividades}
\begin{enumerate}
  \item Organice la serie en Excel y represente gráficamente población
    contra tiempo.
  \item Reordene la serie si lo considera necesario para analizarla
    cronológicamente y describa la tendencia observada.
  \item Calcule tasas o cambios medios para subperíodos relevantes e
    indique con claridad si están expresados por año o por período
    censal.
  \item Aplique los siguientes métodos:
    \begin{itemize}
      \item aritmético,
      \item geométrico,
      \item un ajuste por mínimos cuadrados lineal,
      \item un ajuste por mínimos cuadrados exponencial,
      \item una discusión exploratoria del ajuste logístico.
    \end{itemize}
  \item Proyecte la población al año \textbf{2046}.
  \item Compare los resultados de los métodos y discuta cuál describe
    mejor el comportamiento histórico implícito en la serie.
  \item Seleccione una población de proyecto razonable y justifique su
    elección.
\end{enumerate}

\section{Parte C. Ejercicio con datos reales: Restauración, Dajabón}

\subsection{Contexto}
En esta parte se trabaja con una serie histórica real, más extensa que
la de La Siembra. El propósito es contrastar un ejercicio controlado con
un caso censal real, donde el comportamiento no necesariamente sigue un
modelo simple.

\subsection{Datos}
Los valores siguientes deben ser verificados y citados contra la ONE en
la entrega final.

\begin{center}
\begin{tabular}{l *{7}{S[table-format=4.0]}}
  \toprule
  {Año censal} & {1960} & {1970} & {1981} & {1993} & {2002} & {2010} & {2022} \\
  \midrule
  {Población (hab)} & 7259 & 7243 & 7261 & 6954 & 6187 & 6252 & 5728 \\
  \bottomrule
\end{tabular}
\end{center}

\subsection{Actividades}
\begin{enumerate}
  \item Organice la serie en Excel y represente gráficamente la
    evolución poblacional.
  \item Analice el comportamiento de la serie, identifique cambios de
    tendencia y comente posibles implicaciones para diseño.
  \item Calcule incrementos absolutos, incrementos porcentuales y tasas
    representativas. Si emplea promedios entre intervalos censales
    desiguales, justifique el criterio usado.
  \item Aplique los siguientes métodos:
    \begin{itemize}
      \item aritmético,
      \item geométrico,
      \item ajuste lineal por mínimos cuadrados,
      \item ajuste exponencial por mínimos cuadrados,
      \item un método adicional de su elección entre los discutidos en
        clase.
    \end{itemize}
  \item Proyecte la población al año \textbf{2046}.
  \item Compare resultados en tablas y gráficos.
  \item Seleccione una población de proyecto y justifique su elección
    conforme al comportamiento histórico y al Reglamento del INAPA.
\end{enumerate}

\section{Entregables}

\subsection{Archivos requeridos}
Cada grupo debe entregar:
\begin{enumerate}
  \item un archivo Excel con datos, cálculos, fórmulas visibles y
    gráficos;
  \item un informe breve en PDF, de aproximadamente 4 a 6 páginas de
    contenido, que sintetice procedimiento, resultados, comparación de
    métodos y población de proyecto adoptada.
\end{enumerate}

\subsection{Formato obligatorio de nombres}
Use exactamente la siguiente estructura:
\begin{itemize}
  \item \texttt{GXX\_excel\_proy\_pob.xlsx}
  \item \texttt{GXX\_informe\_proy\_pob.pdf}
\end{itemize}
donde \texttt{GXX} identifica al grupo, por ejemplo, \texttt{G03}.

\begin{icvwarning}[Entrega]
El incumplimiento del formato de nombre de archivo tendrá una
penalización de \textbf{10 puntos} sobre la calificación total de la
asignación.
\end{icvwarning}

\subsection{Contenido mínimo del informe}
El informe deberá incluir, como mínimo:
\begin{itemize}
  \item breve introducción;
  \item metodología utilizada;
  \item resultados de la Parte A;
  \item resultados de la Parte B;
  \item resultados de la Parte C;
  \item comparación de métodos;
  \item justificación final de la población de proyecto;
  \item conclusiones.
\end{itemize}

\section{Pregunta integradora}
Desde el punto de vista del diseño de sistemas de agua potable y
alcantarillado, explique qué riesgos técnicos implica una mala
estimación de la población de proyecto en comunidades pequeñas o con
comportamiento demográfico inestable.

\section{Relación con el proyecto de la asignatura}
La metodología aplicada en esta asignación servirá como base para
justificar población de diseño, gastos de diseño y decisiones
preliminares de dimensionamiento en etapas posteriores del curso.

\appendix
\section{Guía breve de métodos permitidos}

\subsection{Método aritmético}
Supone incremento absoluto aproximadamente constante en el tiempo.
Suele ser razonable cuando el crecimiento es lento o relativamente
estable en el horizonte analizado.

\subsection{Método geométrico}
Supone una tasa porcentual aproximadamente constante. Debe usarse con
cautela en comunidades pequeñas o con señales de estancamiento o
decrecimiento.

\subsection{Ajuste lineal por mínimos cuadrados}
Permite ajustar una recta a los datos históricos. Es útil cuando la
tendencia global es aproximadamente lineal dentro del período
estudiado.

\subsection{Ajuste exponencial por mínimos cuadrados}
Permite ajustar un modelo exponencial cuando los datos sugieren cambios
proporcionales más que incrementos absolutos constantes. Su pertinencia
debe ser justificada.

\subsection{Ajuste logístico}
Puede explorarse cuando la serie sugiere una trayectoria de tipo S y un
límite práctico de crecimiento. En esta asignación interesa más el
razonamiento técnico que la exactitud de un ajuste perfecto.

\section{Rúbrica de evaluación}

\begin{center}
\begin{tabular}{p{0.28\linewidth} p{0.52\linewidth} p{0.10\linewidth}}
  \toprule
  \textbf{Criterio} & \textbf{Descripción} & \textbf{Puntos} \\
  \midrule
  Identificación y verificación de datos censales &
    Identifica correctamente los datos, cita la ONE, documenta
    discrepancias cuando existan y distingue datos observados de
    valores proyectados. & 20 \\
  Organización y análisis en Excel &
    Presenta hojas claras, fórmulas consistentes, cálculo trazable,
    unidades bien indicadas y gráficos legibles. & 15 \\
  Aplicación de métodos de proyección &
    Aplica correctamente los métodos requeridos y explicita supuestos,
    alcance y limitaciones de cada uno. & 25 \\
  Análisis crítico y comparación &
    Compara resultados con criterio ingenieril y explica por qué
    algunos métodos pueden ser menos adecuados que otros. & 20 \\
  Justificación de la población de proyecto &
    Selecciona una población de proyecto técnicamente defendible y
    alineada con el Reglamento del INAPA. & 20 \\
  \bottomrule
\end{tabular}
\end{center}

\begin{icvwarning}[Observaciones de evaluación]
\begin{itemize}
  \item No existe una única respuesta correcta.
  \item Se calificará prioritariamente la solidez del razonamiento
    técnico.
  \item El uso de fuentes oficiales de la ONE es obligatorio.
  \item El incumplimiento del formato de nombres de archivo descuenta
    \textbf{10 puntos} de la nota final.
\end{itemize}
\end{icvwarning}

\section{Checklist de revisión rápida}

\begin{longtable}{p{0.95\linewidth}}
  \toprule
  \textbf{Fuentes de información y trazabilidad} \\
  \midrule
  $\square$ Se identifica claramente el año censal de cada dato
    utilizado. \\
  $\square$ Se cita explícitamente la ONE como fuente de los datos
    verificados. \\
  $\square$ Se distingue entre dato observado y valor proyectado. \\
  $\square$ En La Siembra, se documentan las limitaciones o
    discrepancias encontradas. \\
  \midrule
  \textbf{Organización y claridad de los datos} \\
  \midrule
  $\square$ Los datos están organizados de forma clara en Excel. \\
  $\square$ Las unidades y los años están correctamente indicados. \\
  $\square$ Las tablas y los gráficos son coherentes entre sí. \\
  $\square$ No hay inconsistencias numéricas evidentes. \\
  \midrule
  \textbf{Métodos de proyección aplicados} \\
  \midrule
  $\square$ Se aplican los métodos requeridos en cada apartado. \\
  $\square$ Las fórmulas utilizadas son correctas. \\
  $\square$ Los supuestos de cada método están explícitos. \\
  $\square$ Se reconoce cuándo un método puede ser poco adecuado para
    el caso analizado. \\
  \midrule
  \textbf{Análisis crítico y justificación} \\
  \midrule
  $\square$ Se comparan resultados obtenidos con distintos métodos. \\
  $\square$ Se discuten las limitaciones de los datos disponibles. \\
  $\square$ Se selecciona una población de proyecto y se justifica
    técnicamente. \\
  $\square$ La justificación es coherente con el Reglamento del
    INAPA. \\
  \midrule
  \textbf{Presentación y entrega} \\
  \midrule
  $\square$ El informe es claro, ordenado y conciso. \\
  $\square$ Los archivos entregados siguen el formato obligatorio de
    nombres. \\
  $\square$ Las conclusiones responden al objetivo de la actividad. \\
  $\square$ El trabajo demuestra criterio ingenieril. \\
  \bottomrule
\end{longtable}

\end{document}
